% Exemplo: prova com várias versões. As alternativas são embaralhadas pelo pacote
% exam-randomizechoices, com o número da versão como semente; o gabarito sai do mesmo arquivo.
% Compilar a versão 1 (padrão):  latexmk -pdf prova-versoes.tex
% Versão 3:          latexmk -pdf -jobname=prova-v3 -usepretex="\def\versao{3}" prova-versoes.tex
% Gabarito da 3:     latexmk -pdf -jobname=gabarito-v3 -usepretex="\def\versao{3}\def\gabarito{1}" prova-versoes.tex
% Várias versões de uma vez: scripts gerar-versoes.sh (Linux/macOS) e gerar-versoes.ps1 (Windows).
% Testado com: exam 2.7, exam-randomizechoices, TeX Live 2024 ou mais recente
\documentclass[12pt, a4paper, addpoints]{exam}
\usepackage[T1]{fontenc}
\usepackage[brazilian]{babel}
\usepackage{amssymb}
\pointpoints{ponto}{pontos}
\renewcommand{\solutiontitle}{\noindent\textbf{Resposta:}\enspace}
\CorrectChoiceEmphasis{\bfseries}
\renewcommand{\thechoice}{\alph{choice}}     % alternativas a), b), c)...
\renewcommand{\choicelabel}{\thechoice)}

%%BEGIN:versoespreambulo
\usepackage{exam-randomizechoices}   % ambiente randomizechoices: alternativas em ordem sorteada
\providecommand{\versao}{1}          % versão padrão, se nenhuma for informada na compilação
\setrandomizerseed{\versao}          % a versão é a semente do sorteio: mesma versão, mesma ordem
\ifdefined\gabarito \printanswers \fi
%%END:versoespreambulo
\keylistquestionname{Questão}\keylistkeyname{Resposta}
% letra da alternativa correta na versão compilada (acompanha o sorteio)
\newcommand{\altcorreta}{\ref{question@\thequestion @correctchoice}}

\pagestyle{headandfoot}
\firstpageheader{Introdução ao \LaTeX}{Prova 1 --- versão \versao}{\today}
\runningheader{Introdução ao \LaTeX}{Prova 1 --- versão \versao}{}
\footer{}{Página \thepage\ de \numpages}{}

%%BEGIN:folharespostas
% Folha de respostas: os números das questões vêm dos \label de cada questão
\newcommand{\circulo}[1]{\textcircled{\footnotesize #1}}
\newcommand{\linhaMC}[1]{\ref{#1} & \circulo{a} & \circulo{b} & \circulo{c} & \circulo{d}\\[1ex]}
\newcommand{\folharespostas}{%
  \begin{center}
    \fbox{\begin{tabular}{c|cccc}
      \multicolumn{5}{c}{\textbf{Folha de respostas: marque uma alternativa}}\\[1ex]
      \textbf{Questão} & a & b & c & d \\ \hline\rule{0pt}{2.6ex}
      \linhaMC{q:comentario}
      \linhaMC{q:biber}
      \linhaMC{q:label}
    \end{tabular}}
  \end{center}}
%%END:folharespostas

\begin{document}

\noindent Nome: \makebox[0.55\textwidth]{\hrulefill} \hfill Matrícula: \makebox[3cm]{\hrulefill}

\folharespostas

%%BEGIN:versoesquestoes
\begin{questions}
  \question[1]\label{q:comentario}
  Qual caractere inicia um comentário em um arquivo \texttt{.tex}?
  \begin{randomizechoices}
    \choice \texttt{\#}
    \correctchoice \texttt{\%}
    \choice \texttt{\&}
    \choice \texttt{\textbackslash}
  \end{randomizechoices}
  \begin{solution}
    Alternativa \altcorreta: o \texttt{\%} faz o \LaTeX\ ignorar o resto da linha.
  \end{solution}

  \question[1]\label{q:biber}
  Com o pacote \texttt{biblatex}, qual programa monta a lista de referências?
  \begin{randomizechoices}
    \choice \texttt{makeindex}
    \choice \texttt{bibtex8}
    \correctchoice \texttt{biber}
    \choice \texttt{dvips}
  \end{randomizechoices}
  \begin{solution}
    Alternativa \altcorreta.
  \end{solution}

  \question[1]\label{q:label}
  Em uma figura, onde deve ficar o \texttt{\textbackslash label}?
  \begin{randomizechoices}
    \choice Antes do \texttt{\textbackslash begin\{figure\}}.
    \choice Antes do \texttt{\textbackslash caption}.
    \correctchoice Logo depois do \texttt{\textbackslash caption}.
    \choice Depois do \texttt{\textbackslash end\{figure\}}.
  \end{randomizechoices}
  \begin{solution}
    Alternativa \altcorreta: o rótulo guarda o número criado pelo \texttt{\textbackslash caption}.
  \end{solution}

  \question[2]
  Explique a diferença entre uma tabela e um quadro nos trabalhos acadêmicos.
  \begin{solutionorlines}[2.5cm]
    A tabela apresenta dados, geralmente numéricos, com as laterais abertas; o quadro apresenta
    informação textual, com bordas em todos os lados.
  \end{solutionorlines}
\end{questions}
%%END:versoesquestoes

%%BEGIN:versoeschave
\ifprintanswers
  \bigskip\noindent\textbf{Chave de respostas --- versão \versao}\par\smallskip
  \printkeytable[1-3]
\fi
%%END:versoeschave

\end{document}
